%%%%%%%%%%%%%%%%%%%%%%%%%%%%%%%%%%%%%%%%%%%%%%%%%%%%%%%%%%%%%%%%%%%%%%
% Lớp 8 | Chương III. Tứ giác | Hình chữ nhật
% Dạng 8H3?13-1: Tính chất hình chữ nhật, áp dụng vào tam giác vuông
%%%%%%%%%%%%%%%%%%%%%%%%%%%%%%%%%%%%%%%%%%%%%%%%%%%%%%%%%%%%%%%%%%%%%%


% nguồn: De_AI_Toan8_cuoi_hoc_ki_i_1 (nhập 21/11/2026)
\begin{bt}%[8H3H13-1]%[ma:897169ff]
	Cho hình chữ nhật $ABCD$ có $AB = 8\text{ cm}$, $BC = 6\text{ cm}$. Gọi $O$ là giao điểm của hai đường chéo $AC$ và $BD$.
a) Tính độ dài các đường chéo $AC$ và $BD$.
b) Tính độ dài các đoạn thẳng $OA$, $OB$, $OC$, $OD$.
	\loigiai{
		a) Xét tam giác $ABC$ vuông tại $B$, theo định lí Pythagore ta có: \\
$AC^2 = AB^2 + BC^2 = 8^2 + 6^2 = 64 + 36 = 100 \Rightarrow AC = 10\text{ cm}$. \\
Vì $ABCD$ là hình chữ nhật nên hai đường chéo bằng nhau: $BD = AC = 10\text{ cm}$. \\
b) Trong hình chữ nhật, hai đường chéo bằng nhau và cắt nhau tại trung điểm của mỗi đường, nên: \\
$OA = OB = OC = OD = \dfrac{AC}{2} = \dfrac{10}{2} = 5\text{ cm}$.
	}
\end{bt}

% nguồn: De_AI_Toan8_cuoi_hoc_ki_i_1 (nhập 21/11/2026)
\begin{bt}%[8H3H13-1]%[ma:8f4bbf4e]
	Cho tam giác $ABC$ vuông tại $A$ có $AB = 9\text{ cm}$, $AC = 12\text{ cm}$. Gọi $M$ là trung điểm của cạnh huyền $BC$.
a) Tính độ dài cạnh huyền $BC$.
b) Tính độ dài đường trung tuyến $AM$.
	\loigiai{
		a) Áp dụng định lí Pythagore trong tam giác $ABC$ vuông tại $A$: \\
$BC^2 = AB^2 + AC^2 = 9^2 + 12^2 = 81 + 144 = 225 \Rightarrow BC = 15\text{ cm}$. \\
b) Trong tam giác vuông, đường trung tuyến ứng với cạnh huyền bằng nửa cạnh huyền, nên: \\
$AM = \dfrac{1}{2}BC = \dfrac{15}{2} = 7{,}5\text{ cm}$.
	}
\end{bt}

% nguồn: De_AI_Toan8_cuoi_hoc_ki_i_1 (nhập 21/11/2026)
\begin{bt}%[8H3H13-1]%[ma:01050acb]
	Cho hình chữ nhật $ABCD$ có $O$ là giao điểm của hai đường chéo $AC$ và $BD$. Biết $AB = 5\text{ cm}$ và $\widehat{AOB} = 60^\circ$.
a) Chứng minh tam giác $OAB$ là tam giác đều.
b) Tính độ dài đường chéo $AC$ và cạnh $BC$.
	\loigiai{
		a) Trong hình chữ nhật $ABCD$, hai đường chéo bằng nhau và cắt nhau tại trung điểm $O$, nên $OA = OB$. \\
Suy ra tam giác $OAB$ cân tại $O$. Lại có $\widehat{AOB} = 60^\circ$ nên tam giác $OAB$ là tam giác đều. \\
b) Vì tam giác $OAB$ đều nên $OA = OB = AB = 5\text{ cm}$. \\
Suy ra $AC = 2OA = 2 \cdot 5 = 10\text{ cm}$. \\
Áp dụng định lí Pythagore trong tam giác $ABC$ vuông tại $B$: $BC^2 = AC^2 - AB^2 = 10^2 - 5^2 = 75 \Rightarrow BC = \sqrt{75} = 5\sqrt{3}\text{ cm}$.
	}
\end{bt}

% nguồn: De_AI_Toan8_cuoi_hoc_ki_i_1 (nhập 21/11/2026)
\begin{bt}%[8H3H13-1]%[ma:3d636a60]
	Cho tam giác $ABC$ có đường trung tuyến $AM$ thỏa mãn $AM = \dfrac{1}{2}BC$.
a) Chứng minh tam giác $ABC$ vuông tại $A$.
b) Biết $AB = 7\text{ cm}$ và $BC = 25\text{ cm}$, tính độ dài cạnh $AC$.
	\loigiai{
		a) Vì $M$ là trung điểm của $BC$ và $AM = \dfrac{1}{2}BC$ nên $MA = MB = MC$. \\
Tam giác $MAB$ cân tại $M$ nên $\widehat{MAB} = \widehat{B}$; tam giác $MAC$ cân tại $M$ nên $\widehat{MAC} = \widehat{C}$. \\
Trong tam giác $ABC$, ta có $\widehat{A} + \widehat{B} + \widehat{C} = 180^\circ \Rightarrow \widehat{A} + (\widehat{MAB} + \widehat{MAC}) = 180^\circ \Rightarrow 2\widehat{A} = 180^\circ \Rightarrow \widehat{A} = 90^\circ$. \\
Vậy tam giác $ABC$ vuông tại $A$. \\
b) Áp dụng định lí Pythagore trong tam giác $ABC$ vuông tại $A$: $AC^2 = BC^2 - AB^2 = 25^2 - 7^2 = 625 - 49 = 576 \Rightarrow AC = 24\text{ cm}$.
	}
\end{bt}

% nguồn: chuyen_dang (nhập 22/11/2026)
\begin{bt}%[8H3C13-1]%[ma:c1a17409]
	Cho tam giác $ABC$ vuông tại $A$ có $\tan B=\dfrac34$ và $AB+AC=14\text{ cm}$. Tính độ dài cạnh $BC$ (đơn vị: cm).
	\shortans{10}
	\loigiai{
		$\tan B=\dfrac{AC}{AB}=\dfrac34$ nên đặt $AC=3k$, $AB=4k$ ($k>0$).\\
		$AB+AC=7k=14\Rightarrow k=2$, suy ra $AB=8\text{ cm}$, $AC=6\text{ cm}$.\\
		$BC=\sqrt{8^2+6^2}=10\text{ (cm)}$.
	}
\end{bt}
